% Diagram konsep Bab 13: Indikator yang berbeda memicu keputusan yang berbeda.
% Dibangkitkan tools/render-diagrams.py dari figures/chapter-diagrams.json;
% ubah data JSON lalu jalankan ulang skrip, jangan sunting berkas ini langsung.
\begin{tikzpicture}[
  x=0.92cm, y=0.95cm,
  kotak konsep/.style={draw, line width=0.5pt, rounded corners=2pt,
    align=center, minimum width=2.999cm,
    minimum height=1.330cm, inner sep=0pt, font=\scriptsize,
    text=konsepJudul},
  peran actor/.style={kotak konsep, fill=konsepIsiActor, draw=konsepGarisActor},
  peran service/.style={kotak konsep, fill=konsepIsiService, draw=konsepGarisService},
  peran control/.style={kotak konsep, fill=konsepIsiControl, draw=konsepGarisControl},
  peran data/.style={kotak konsep, fill=konsepIsiData, draw=konsepGarisData},
  peran outcome/.style={kotak konsep, fill=konsepIsiOutcome, draw=konsepGarisOutcome},
  peran risk/.style={kotak konsep, fill=konsepIsiRisk, draw=konsepGarisRisk},
  panah konsep/.style={-{Stealth[length=1.3mm, width=1.2mm]}, line width=0.6pt,
    draw=konsepPanah},
  panah putus/.style={panah konsep, draw=konsepPanahPutus,
    dash pattern=on 2.5pt off 1.5pt},
  label panah/.style={font=\tiny, align=center, fill=white,
    inner sep=1pt, text=konsepLabel},
  catatan/.style={font=\scriptsize, align=left, anchor=north west,
    inner sep=0pt, text width=12.383cm, text=konsepTeks,
    execute at begin node={\hyphenpenalty=10000\hbadness=10000\relax}},
]
  \definecolor{konsepIsiActor}{HTML}{EEF1F6}
  \definecolor{konsepGarisActor}{HTML}{50627D}
  \definecolor{konsepIsiService}{HTML}{EAF2FB}
  \definecolor{konsepGarisService}{HTML}{3F6893}
  \definecolor{konsepIsiControl}{HTML}{E9F3EC}
  \definecolor{konsepGarisControl}{HTML}{316B4B}
  \definecolor{konsepIsiData}{HTML}{F1EDF8}
  \definecolor{konsepGarisData}{HTML}{786394}
  \definecolor{konsepIsiOutcome}{HTML}{E4F2EC}
  \definecolor{konsepGarisOutcome}{HTML}{236C57}
  \definecolor{konsepIsiRisk}{HTML}{FCEDE7}
  \definecolor{konsepGarisRisk}{HTML}{B36947}
  \definecolor{konsepJudul}{HTML}{20394F}
  \definecolor{konsepTeks}{HTML}{455264}
  \definecolor{konsepLabel}{HTML}{35485C}
  \definecolor{konsepPanah}{HTML}{5C7084}
  \definecolor{konsepPanahPutus}{HTML}{A15D40}
  % Panah dan label lebih dahulu; kotak menutupi sisa latar label.
  \draw[panah konsep] (0.93,0) -- (2.77,0);
  \draw[panah konsep] (-0.7,-0.7) -- (-0.7,-2);
  \draw[panah konsep] (-0.7,-0.7) -- (-0.7,-0.9) -- (4.4,-0.9) -- (4.4,-2);
  \draw[panah konsep] (6.03,0) -- (7.87,0);
  \draw[panah konsep] (-0.7,-3.4) -- (-0.7,-3.6) -- (9.5,-3.6) -- (9.5,-3.4);
  \draw[panah konsep] (6.03,-2.7) -- (7.87,-2.7);
  \draw[panah konsep] (9.5,-2) -- (9.5,-0.7);
  \draw[panah putus] (0.93,-2.7) -- (2.77,-2.7);
  \node[label panah] at (1.85,0.31) {hitung hasil};
  \node[label panah] at (-0.7,-1.35) {hitung pelaksanaan};
  \node[label panah] at (4.4,-1.45) {hitung paparan};
  \node[label panah] at (6.95,0.31) {status hasil};
  \node[label panah] at (4.4,-3.91) {bandingkan batas};
  \node[label panah] at (6.95,-2.26) {bandingkan\\batas};
  \node[label panah] at (9.5,-1.35) {usulan tindakan};
  \node[label panah] at (1.85,-2.39) {memengaruhi};
  \node[peran data] (data) at (-0.7,0)
    {\textbf{Data bersumber}\\[0.5pt]{\tiny\color{konsepTeks}populasi, periode, definisi}};
  \node[peran outcome] (kpi) at (4.4,0)
    {\textbf{KPI hasil}\\[0.5pt]{\tiny\color{konsepTeks}tujuan layanan tercapai?}};
  \node[peran control] (putusan) at (9.5,0)
    {\textbf{Keputusan direksi}\\[0.5pt]{\tiny\color{konsepTeks}opsi, biaya, pemilik, tenggat}};
  \node[peran control] (kci) at (-0.7,-2.7)
    {\textbf{KCI kontrol}\\[0.5pt]{\tiny\color{konsepTeks}tinjauan akses terlaksana?}};
  \node[peran risk] (kri) at (4.4,-2.7)
    {\textbf{KRI paparan}\\[0.5pt]{\tiny\color{konsepTeks}akun tanpa pemilik?}};
  \node[peran control] (ambang) at (9.5,-2.7)
    {\textbf{Ambang dan eskalasi}\\[0.5pt]{\tiny\color{konsepTeks}tindakan menurut toleransi}};
  \node[catatan] at (-2.33,-4.5)
    {KCI tinjauan 95\% dapat hijau sementara satu akun tanpa pemilik membuat KRI merah. Kedua status tidak dirata-ratakan.};
\end{tikzpicture}
